\begin{table}[tbp]
    \centering
    \small
    \caption{Cronograma executivo de 6 meses para finalização e defesa da dissertação de mestrado.}
    \label{tab:cronograma}
    \begin{tabular}{clp{11.5cm}}
        \toprule
        Mês & Período & Entregáveis e Metas Metodológicas                                                                                                                                             \\
        \midrule
        M1  & Mês 1   & Fechamento dos capítulos metodológicos da bancada e calibração da rede recorrente CA3 no Brian 2 (Izhikevich: piramidais RS + interneurônios FS PV$^+$).                      \\
        M2  & Mês 2   & Simulação computacional sistemática do protocolo óptico (Laura Costa, 2026: ChR2 e cinética lenta) vs. elétrico (Wagner Franklin: recrutamento inibitório e regime contínuo). \\
        M3  & Mês 3   & Varredura paramétrica de regimes temporais (efeito de \emph{bursts}, \emph{duty cycle} e teste de estimulação em malha fechada / \emph{closed-loop}).                         \\
        M4  & Mês 4   & Redação integral dos Capítulos de Resultados da Simulação e Discussão sob a ótica de Processamento de Sinais em Redes Neurais.                                                \\
        M5  & Mês 5   & Consolidação do texto completo da dissertação e envio aos orientadores (Prof. Márcio Flávio e Prof. Hani Yehia) para correções.                                               \\
        M6  & Mês 6   & Depósito formal junto ao colegiado do PPGEE/UFMG e Defesa da Dissertação de Mestrado.                                                                                         \\
        \bottomrule
    \end{tabular}
\end{table}